\documentclass[11pt]{article}
\usepackage[margin=1in]{geometry}
\usepackage{amsmath,amssymb}
\usepackage{enumitem}
\usepackage{booktabs}

\newcommand{\indep}{{\bot\negthickspace\negthickspace\bot}}
\DeclareMathOperator{\E}{\mathbb{E}}

\pagestyle{plain}

\begin{document}

\begin{center}
{\Large \textbf{PS200C: Quantitative Methods}} \\[6pt]
{\large In-class Activity 5: Estimation under Selection on Observables}
\end{center}

\vspace{4pt}
\noindent\textbf{Name:} \hrulefill

\vspace{12pt}

%%% QUESTION 1 %%%
\noindent\textbf{Question 1: All roads lead to the same place}

\medskip
\noindent Under CI, every method we discussed is estimating
$$\E_{X \sim F}\!\left[\mu_1(X) - \mu_0(X)\right], \qquad \mu_d(x) \equiv \E[Y(d) \mid X = x].$$

\begin{enumerate}[label=(\alph*)]
\item Which distribution $F$ corresponds to the ATE? The ATT? The ATC?

\vspace{2.2cm}

\item What is the key difference between any estimand of this kind and the simple difference in means? For a given choice of $F$, how does this differ from what you would get by comparing a group average for the treated vs. control? Why is this what we want to do when we are ``adjusting for X''?


\vspace{2cm}
\end{enumerate}


%%% QUESTION 2 %%%
\noindent\textbf{Question 2: IPW by hand}

\medskip
\noindent Recall that the (fitted) propensity score, $\hat{\pi}(X_i)$ is an estimate of the probability unit $i$ takes the treatment, based on its $X_i$, i.e. $Pr(D_i=1 \mid X_i=1)$. Four units, each with an estimated propensity score $\widehat{\pi}(X_i)$:

\begin{center}
\begin{tabular}{@{} c c c c @{}}
\toprule
Unit & $D_i$ & $\widehat{\pi}(X_i)$ & $Y_i$ \\
\midrule
1 & 1 & 0.80 & 10 \\
2 & 1 & 0.20 & \phantom{0}6 \\
3 & 0 & 0.80 & \phantom{0}8 \\
4 & 0 & 0.20 & \phantom{0}4 \\
\bottomrule
\end{tabular}
\end{center}

\begin{enumerate}[label=(\alph*)]
\item Compute the simple IPW weight $w_i = D_i/\widehat{\pi}(X_i) + (1-D_i)/(1-\widehat{\pi}(X_i))$ for each unit.

\vspace{2.2cm}

\item Which unit gets upweighted the most, and why does that make intuitive sense?

\vspace{1.8cm}
\end{enumerate}

\newpage

%%% QUESTION 3 %%%
\noindent\textbf{Question 3: Matching the method to the problem}

\medskip
\noindent For each scenario, name one estimation approach you'd prefer and one you'd avoid, with a one-line reason. Some choices include sub-classification/stratification, matching on $X$, matching on the pscore, IPW, balancing weights, and regression.

\begin{enumerate}[label=(\alph*)]

\item Just 3 discrete covariates (sex, age bin, party ID), $n = 5000$, with plenty of both treated and control in every cell.

\vspace{1.8cm}
\item You have 10 covariates and $n = 300$.  Treated units are rare in some covariate combinations. However, you suspect (or hope) that not all variables really influence $Pr(D)$ much.



\vspace{1.8cm}
\end{enumerate}


%%% QUESTION 4 %%%
\noindent\textbf{Question 4: Regression as imputation}

\medskip
\noindent In the g-computation / regression-imputation view of regression:

\begin{enumerate}[label=(\alph*)]
\item What do you fit $\widehat{\mu}_1(x)$ on, and what does it estimate?  Same question for $\widehat{\mu}_0(x)$.


\vspace{2cm}

\item When treatment effects vary with $X$, why is fitting \emph{separate} linear models on treated and control (or equivalently, the Lin regression) preferable to pooled OLS with $Y \sim D + X$?

\vspace{2cm}

\end{enumerate}


%%% QUESTION 5 %%%
\noindent\textbf{Question 5: The big picture}

\medskip
\noindent Does using a more sophisticated adjustment technique, and/or having a larger sample, help reduce concerns about omitted confounders?  Why or why not?

\vspace{2.5cm}

\end{document}
